\documentclass[twoside]{article}
\usepackage[
  a5paper,
  margin=0.5in,
  includehead,
  headsep=4pt,
  headheight=13pt
]{geometry}
\usepackage{gregosheet}
\usepackage{gregosheet-psalm}
\usepackage{lyluatex}
\usepackage{xcolor}
\input{styles}

\begin{document}

\section*{Temetés szertartása}

\parttitle[TSZ 24]{Kezdőének}
\begin{gregosheet}
  \addinline{<X-3--3---4---3---4--5-:-4-5--6---5---4--3-,-4--3-2----1-0--3--4---3-:-4-5---6---5---4--3-.}
            {   Kö-rül vet-tek en-gem a ha-lál kín-ja-i;  Ó, U-ram, a te ke-zed-be  a-ján-lom lel-ke-met.}
\end{gregosheet}
%\begin{gregosheet}
%  \addinline{<X-3--3---4---3---4--5-:-4-5--6---5---4--3-,-4--3-2----1-0--3--4---3-:-4-5---6---5---4--3-?,-5--7-7---7--5----6---6--5--:-5----[---------4---5---6---5----4---3-,-5--¦-------------5----6---5--:-¦------------4---5--6--5--4---3-?,-3--3--?,--5-¦----------------¦------------------------¨-:-¦----------------------------5--6--6--5--:-5---5--5----4------5---6--5--4---3--3--,--5----¦--------------5--6-5--:-5--4---5---6-5--4--3--3--?,-3--3--?,--5--7---7----7---7---7---5---6--5--:-5----5--5--4---5---6--5---4--3-,-5----7---7---¦--------------5-6--6---5--:-[-----------------------------------4---5--6--5---4---3-?,-3--3--?,--5--¦--------------------5----6--5-:-5--4-5---6--5---4--3--,-5----¦------------------5---6--5-:-5--5--[---------------4--5----6---5---4--3-?,-3--3-.}
%            {   Kö-rül vet-tek en-gem a ha-lál kín-ja-i;  Ó, U-ram, a te ke-zed-be  a-ján-lom lel-ke-met. Az U-rat na-gyon sze-re-tem, mert imádságom sza-vát meg-hall-gat-ta. És kegyesen_leha-jolt hoz-zám, valahányszor hoz-zá ki-ál-tot-tam. Kö-rül... A keserűség_és_fáj-dalom_szorongatot_engem, @   én_azonban_segítségül_hívtam az Úr ne-vét: „Is-te-nem, mentsd meg az én lel-ke-met!” Mert irgalmas_az_úr és i-gaz, kö-nyö-rül a mi Is-te-nünk. Kö-rül... Én lel-kem, nyu-god-jál meg im-már, mert az Úr gon-dot vi-sel re-ád! Mert meg-sza-badított_engem a ha-lál-tól, szememet_a_könnyhullatástól_és_lába-mat az el-bot-lás-tól. Kö-rül... És így_az_Úrnak_színe_e-lőtt já-rok az é-lők or-szá-gá-ban. Adj, Uram_örök_nyugodal-mat ne-ki, és az örök_világosság fé-nyes-ked-jék ne-ki!  Kö-rül...}
%\end{gregosheet}

\parttitle[TSZ 27]{Válaszos zsoltár}
\begin{gregosheet}
  \addinline{<ô-Ÿ-------------2--1---1r-4-:-£-------------2---3r--5----,-£------------------2----1-2r--4-.}
            {   Az_Úr_irgalma kí-sér en-gem életemnek_min-den nap-ján, * az_ő_házában_lakom mind-ö-rök-ké.}
\end{gregosheet}
%V. Az Úr nékem pászt\underline{o}rom, † nincs hiányom semmi j\underline{ó}ban, * ő vezet az igazs\underline{á}g \underline{ú}tján. \textbf{AZ Ő HÁZÁBAN...} \\
%V. Zöldellő réteken vezet \underline{e}ngem, † fölnevel szelíd vizek m\underline{e}ntén, * fölüdíti fár\underline{a}dt l\underline{e}lkem. \textbf{AZ Ő HÁZÁBAN...} \\
%V. Járjam bár a halál v\underline{ö}lgyét, † nem félek semmi b\underline{a}jtól, * mert te ott v\underline{a}gy v\underline{é}lem. \textbf{AZ Ő HÁZÁBAN...}

\parttitle[TSZ 33]{Búcsúvételi ének}
\begin{gregosheet}
  \addinline{<-2---4--5---[-----------------rf2---4--5-,-[----------------4--2---4--5-?,-2--4---[--------5z-7---5--,-[---------------------4---2--4-5-.}
            {  Jöj-je-tek segítségül,_Isten szent-je-i!  Siessetek,_Urunk an-gya-la-i! * Fo-gad-játok_az ő  lel-két, s_vigyétek_a_Fölséges szí-ne e-lé!}
\end{gregosheet}
%\begin{gregosheet}
%  \addinline{<-¨-----2---4--5---[-----------------rf2---4--5-,-[----------------4--2---4--5-?,-¨-----2--4---5--5---5--5z-7---5--,-[---------------------4---2--4-5-?,-¨-----¡-----------------------------3---2---1---2-,-0--1---3--3---3--:-¢------------------------4-5--4--2-?,-¨-----2--4---5--5--?,--¨-----¡--------------------3---2--1---2--,-0--1--¢----------------------4---5---4--2-?,-¨-----2--4---5--5--?,--¨-----2---4--5--?,--¨-----2--4---5--5-.}
%            {  <E.:> Jöj-je-tek segítségül,_Isten szent-je-i!  Siessetek,_Urunk an-gya-la-i! * <H.:> Fo-gad-já-tok az ő  lel-két, s_vigyétek_a_Fölséges szí-ne e-lé!  <E.:> Fogadd_magadhoz,_Krisztus,_ki meg-hív-tad őt, és ti, an-gya-lok, vezessétek_a_mennyország ö-rö-mé-be!  <H.:> Fo-gad-já-tok... <E.:> Adj,_Uram,_neki_örök nyu-go-dal-mat, és az örök_világosság_fényes-ked-jék ne-ki!  <H.:> Fo-gad-já-tok... <E.:> Jöj-je-tek... <H.:> Fo-gad-já-tok...}
%\end{gregosheet}

\parttitle{Gyászmenet énekei}
\begin{gregosheet}
  \addinline{<X-4-----4--4----2-3--4--4-:-4--5---4---5--4--3---2--,-2--3--4---5--4x--1---2----3--4---3-:-2-3--4---5--4x---3----2-1---0í-.}
            {   Menny-or-szág ö-rö-mé-be  vi-gye-nek az an-gya-lok, és az egy-ko-ron sze-gény Lá-zár-ral é-le-ted le-gyen mind-ö-rök-ké!}
\end{gregosheet}

\textcolor{red}{9. tónus:} Vedd fontolóra, Urunk, h\underline{o}gy porból lettünk, * és elhervad az ember, mint a fűszál, vagy \underline{a} rét virága!

\textcolor{red}{5. tónus:} Testbe öltözöm \underline{ú}jra, * és látom üdvözítő \underline{I}stenem arcát.

\textcolor{red}{4. a tónus:} \e Lázár oszladozó testét feltámasztott\underline{a}d a sírból. * \h Adj halottunknak nyugodalmat, és \underline{e}ngedd el bűneit!

\textcolor{red}{11. tónus:} \e ... Vitted világosságodat, hogy meglássanak t\underline{é}ged, * \h Akiket a sötétség börtön\underline{e} fogva tartott.

\begin{hymnstanzas}
  1. Jézus, világ Megváltója,\\
  üdvözlégy, élet adója,\\
  megfeszített Istenfia,\\
  szent kereszted szívem hívja.\\
  Jézus, add, hogy hozzád térjek,\\
  veled haljak, veled éljek!
\stanzabreak
  2. Jézusom, ha jön a végnap,\\
  ismerj engem magadénak!\\
  Homlokomon piros véred,\\
  tiéd vagyok, úgy ítélj meg!\\
  Jézus, add, hogy hozzád térjek,\\
  veled haljak, veled éljek!
\end{hymnstanzas}

\begin{hymnstanzas}
  1. Krisztus feltámadott!\\
  Halál fején tapodott,\\
  égbe utat mutatott.\\
  Tudom, hogy feltámadok.\\
  Feltámadt Krisztus!
\stanzabreak
  2. Krisztus feltámadott!\\
  Szűnjetek már, bánatok!\\
  Ti is nyomán járjatok,\\
  feltámadást várjatok!\\
  Feltámadt Krisztus!
\end{hymnstanzas}

1. Hol vagy, én szerelmes Jézus Krisztusom? / Hol talállak téged, kegyes Megváltóm? / Adj világosságot, mert fényes vagy, / Nálad nélkül szívem igen bágyadt!

2. Téged az én lelkem kíván keresni, / Mert tenálad nélkül el kell kárhozni, / Ó szerelmes Krisztus, fordulj hozzám, / Vigasztald meg szívem, tekints énrám!

3. Ha fel nem talállak, kegyes Krisztusom, / Elkárhozik lelkem, bizonnyal tudom: / Ó siess hát hozzám, ne távozzál, / Ó szerelmes Krisztus, jöjj, vigasztalj!\\

1. Ments meg engem, Uram! az örök haláltól, / Ama rettenetes napon minden bajtól; / Midőn az ég és föld megfogőnak indulni, / S eljössz a világot lángokban ítélni.

2. Örök nyugodalmat adj, ó Uram, néki, / S örök világosság fényeskedjék nékik / Hogy szent trónod körül hozsannát mondhasson, / S téged, boldogítót, örökké áldhasson.

\begin{hymnstanzas}
  1. Adj, Istenem, boldog halált,\\
  Halálomban elnyugovást,\\
  Adj testünknek jó meghalást,\\
  Dicsőséges feltámadást!
\stanzabreak
  2. Földön olyan az életűnk,\\
  Mint a virág: kóró leszünk;\\
  Elváltozunk, elköltözünk,\\
  Koporsó lesz lakóhelyünk.\\
\stanzabreak
  3. Nem hagyhat el a reményünk,\\
  Halhatatlan igét nyertünk:\\
  Örök lakást benned veszünk,\\
  Míg megújul az életünk.\\
\stanzabreak
  4. Mely órában jön el Urunk,\\
  Azt biztosan nem tudhatjuk,\\
  Krisztus azért inti szívünk,\\
  Hogy mindenkor készen legyünk.\\
\stanzabreak
  5. Azért teljes készülettel\\
  Várjuk Krisztust igaz hittel,\\
  Hogy meghaljunk reménységgel,\\
  S beléphessünk bölcs szüzekkel.\\
\stanzabreak
  6. Senki magát el ne bízza,\\
  Hogy álom az élet útja;\\
  Hogyha fejét bűnre adja,\\
  Késő lehet a bánatja.\\
\stanzabreak
  7. Nagy dicsőség lesz azoknak,\\
  Kik mindvégig hűek voltak;\\
  Mint csillagok úgy ragyognak,\\
  Mikor testben feltámadnak.\\
\stanzabreak
  8. Meglátjuk majd Jézusunkat,\\
  Szentek közt Nagyasszonyunkat;\\
  Kik irgalmat gyakoroltak,\\
  Elnyerik bő jutalmukat.\\
\stanzabreak
  9. Visszakapjuk földi testünk,\\
  Olyan lesz majd mint a lelkünk,\\
  Fáradalmat nem érezünk,\\
  Semmi sem lesz nehéz nekünk.\\
\stanzabreak
  10. Fénylenek a jólelkűek.\\
  Krisztus jobbján összegyűlnek,\\
  Szellemszárnnyal repülhetnek,\\
  Szenvedést már nem éreznek.\\
\stanzabreak
  11. Amit minden hívő remél,\\
  Hogy minket is jobbjára kér,\\
  Az angyalok, szentek közé\\
  S felvisz dicső Atyánk elé.\\
\stanzabreak
  12. Legyen áldott a szent Egység,\\
  Atya, Fiú egy Istenség,\\
  Szentlélekkel örök fölség\\
  Mindörökké dicsértessék!
\end{hymnstanzas}

\parttitle[TSZ 40]{A sírhoz közeledve}
\begin{gregosheet}
  \addinline{<ô-4----4--4---2---rR3-2-:-4--4--2--4--5z-5--,-5----6---5---4--rR3-2-:-2--2---2--1---ñ--1w-2-.}
            {   Nyis-sá-tok meg né--kém ka-pu-it az ég-nek, hogy raj-tuk be-lép-ve, há-lát ad-jak az Úr-nak!}
\end{gregosheet}

\parttitle[TSZ 44]{Szenteltvízhintéskor}
\begin{gregosheet}
  \addinline{<ô-4--4--2---rR3--2-:--4-2---4---5-6---5-,-5--6--5----rR3-2--:-2---1--ñ---1w-2-.}
            {   Az Úr-nál nyug-szom ö-rök-kön-ö-rök-ké; őt vá-lasz-tot-tam, ott la-kom ná-la.}
\end{gregosheet}

\parttitle[TSZ 46]{A koporsó sírba helyezésekor}
\begin{gregosheet}
  \addinline{<ô-£-------------2---rR3-2-:-2--4---4---2--4---5z-5--,-5--5-----5----6-5---4-----rR3-2--:-2--1-ñ---1w-2-.}
            {   Földből_alkot-tál en-gem  és tes-tet ad-tál né-kem, tá-massz fel, U-ram s_Meg-vál-tóm, az i-dők vé-gén!}
\end{gregosheet}

\parttitle[TSZ 52]{Befejezésre}
\begin{gregosheet}
  \addinline{<X-0--0--0r--4-:-£-----------------------5--6--tT4-4--,-£----------4-----4---2--3---4----3---1--0--0--,-¢-------------------1e--2--:-1---0---1---3----2-0---0-.}
            {   Az Úr mon-dá: „Én_vagyok_a_feltámadás és az é---let; aki_bennem hisz, még ha meg-halt is, él-ni fog. És_mindaz,_aki_hisz ben-nem, nem hal meg mind-ö-rök-ké.”}
\end{gregosheet}
\end{document}